\section{Rental markets and experimental design}
We use 2020--24 American Community Survey (ACS) Public Use Microdata Sample household records and FirstMover monthly rental advertisements at first appearance \citep{acs,firstmover}. Incomes and prices are expressed in July 2026 dollars using national CPI-U \citep{cpi}. A lower-income household earns at most the survey-weighted median income of all NYC cash renters (\$64,435 in 2024 dollars). A unit must be in the same borough, cost at most 40\% of monthly household income, and meet a bedroom requirement of $\max(0,\lceil h/2\rceil-1)$ capped at eight for household size $h$. The maximum monthly rent at the group cutoff is \$2,286.

The calibrated study draws from 7,193 recent-mover households and from June--August 2026 arrival catalogues. Requiring valid listing fields, checking geography, filtering nominal asks to \$800--20,000 and deduplicating within month left 66,043 NYC listing records from 72,542 raw arrivals. Staten Island had 26, 26 and 20 listings and could not supply a draw of 64. The Bronx, Brooklyn, Manhattan and Queens account for 97.95\% of the intended recent-mover geographic weight. Outcomes are simulated separately within borough-month. A household or listing may recur across draws. The archive includes the construction rules and losses to cleaning.

\paragraph{Two protocols for drawing markets.}
Ordinary markets draw 64 renter instances with replacement using household weights and 64 distinct listings uniformly. In enriched markets, 32 listings come from an affordability band and 32 from above it. The band ends at the survey-weighted 75th percentile of budgets among lower-income donors who can afford at least one listing in the entire borough-month catalogue and meet its bedroom requirement. We then sample 32 lower-income renters who have an eligible unit in the \emph{drawn} inventory and 32 higher-income renters from their ordinary frame, using survey weights within each pool. The enriched sample is a conditional stress test. In ordinary pilot markets 86.96\% of lower-income instances had no option after geographic standardization.

\paragraph{Holding tastes fixed while altering advice.}
Let $b$'s monthly budget be $c_b$, $r_s$ the asking rent, $Q_s$ the mean of standardized physical attributes (log area and bathrooms), $a_{bs}$ spare bedrooms capped at two, and $d_{bs}$ distance in kilometers from the centroid of the renter's area. The observable score with price penalty $w$ is
\[
   v_{bs}(w)=Q_s-wr_s/c_b+0.5a_{bs}-0.5d_{bs}/5 .
\]
Renter utility on eligible pairs is $v_{bs}(w_p)+0.5\xi_s+\sigma\epsilon_{bs}$, with independent standard normal listing and pair shocks. Listings rank by $2.5\min\{I_b/(12r_s),40/12\}$ plus standard normal pair noise, where $I_b$ is annual income. We use deterministic tie breaking. In the nine-point calibration grid $w_p\in\{1,2,4\}$ and $\sigma\in\{0.50,0.75,1.25\}$, the smallest average absolute error across lower/higher rent-burden quartiles was at $(4,0.50)$ (Section~\ref{sec:burden}).

Policy U samples the quota uniformly from eligible listings. Policy $T_w$ returns the highest-scoring eligible listings under $v_{bs}(w)$. Each arm shows $\min(5,d)$ options to a renter with $d$ feasible units, or $\lceil0.4d\rceil$ in the proportional-quota exercise. The higher-income-only interventions hold the latent rankings, listings and lower-income lists fixed; only higher-income recommendations change. A separate sensitivity varies the behavioral model itself. The outcomes follow from visible applicant-proposing DA. In this model all feasible opportunities are preferred to being unmatched. The simulation does not model endogenous rents, search timing or acceptance of a binding contract.

\paragraph{Estimands and independent variation.}
Within each borough-month, sum matched lower-income renters and divide by the sum of all lower-income renters, including those with no option. Average cell rates using common renormalized borough weights and equal month weights. Each seed block contains an entire set of markets; bootstrap resampling keeps that block and its paired interventions intact. A change of one percentage point means one additional placement per hundred lower-income instances under these weights. Our intervals quantify simulation variability conditional on these catalogues, donor pool and ranking assumptions. The burden exercise instead pools matched renters across calibration draws because its target is the fitted distribution among placements.

\paragraph{Research methods and AI assistance.}
The team supplied the intermediation question, public-data pipeline and pilot protocol. OpenAI Codex assisted with formulating and proving extensions, developing the controlled counterfactuals, implementing solvers and simulation replays, analyzing outputs, preparing figures and drafting the text. The underlying housing observations come from the stated public sources. The trained rankers used numerical housing features. The AI-assisted work also includes their implementation, a published benchmark replication, the proofs about displacement and the observed-preference studies using the sources in Section~\ref{sec:field}. Our checks include independently implemented DA and optimization routines, explicit feasible assignments, exhaustive small-market comparisons and archived block-level outcomes. These are computational checks; the study does not rely on observed renter responses.

